\subsection{Regular intersection at one shared vertex}

Let $\Gamma=(V,E)$ be a finite simple graph and let $G$ be a finite group.
For $R\subseteq V$, write $\widehat{\mathcal S}_R^0$ for the actual regular
regional range on the ambient half-edge space: every complementary slice
belongs to the open-region range $\mathcal S_R^0$. All half-edges crossing
$\partial R$ are retained. The result below gives the regular-coordinate
intersection step in the geometry of
\cite[Theorem~5.4]{Schuch2010PEPS}.
% This regular-coordinate result does not complete Theorem 5.4 for arbitrary
% semi-regular representations or arbitrary G-injective physical site maps.

\begin{theorem}[Restriction and gluing of flat configurations]%
    \label{thm:peps_regular_flat_union}
    \lean{TNLean.PEPS.IsRegularRegionHalfEdgeFlat.of_subset,
        TNLean.PEPS.internalEdge_union_of_no_cross,
        TNLean.PEPS.isRegularRegionHalfEdgeFlat_union_iff}
    \leanok
    \uses{def:peps_regional_half_edge_flat}
    Flatness restricts to subregions. Suppose $R\cap S=\{b\}$ and no edge
    joins $R\setminus S$ to $S\setminus R$. Every internal edge of $R\cup S$
    is then internal to $R$ or to $S$, and a half-edge configuration is flat
    on $R\cup S$ if and only if it is flat on both $R$ and $S$.
\end{theorem}
\begin{proof}\leanok
    Restrict a flatness potential to obtain flatness on a subregion.
    Conversely, let $k_R,k_S$ witness flatness on $R,S$. Set
    $t=k_S(b)^{-1}k_R(b)$ and replace $k_S(v)$ by $k_S(v)t$.
    The new potentials agree at $b$, and for each internal edge $e=(v,w)$
    of $S$ one has
    \begin{align}
        (k_S(v)t)^{-1}\alpha_{v,e}
         & =t^{-1}k_S(v)^{-1}\alpha_{v,e}
        =t^{-1}k_S(w)^{-1}\alpha_{w,e}
        =(k_S(w)t)^{-1}\alpha_{w,e}.
        \label{eq:peps_flat_potential_gluing}
    \end{align}
    Thus the potentials define one potential on $R\cup S$.
    An edge internal to neither region would join their disjoint parts,
    contrary to the hypothesis, so every required equality holds.
\end{proof}

\begin{theorem}[Atomic characterization and nesting of regular ranges]%
    \label{thm:peps_regular_range_nested}
    \lean{TNLean.PEPS.mem_regularGlobalRegionRange_iff_fixed_flat,
        TNLean.PEPS.regularGlobalRegionRange_antitone}
    \leanok
    \uses{def:peps_ambient_regular_range,def:peps_ambient_regional_constraints}
    A vector $x$ lies in $\widehat{\mathcal S}_R^0$ exactly when
    $A_vx=x$ for every $v\in R$, $B_ex=x$ for every internal edge $e$ of $R$,
    and $x(\alpha)=0$ whenever $\alpha|_R$ is not flat.
    Consequently $R\subseteq S$ implies
    $\widehat{\mathcal S}_S^0\subseteq\widehat{\mathcal S}_R^0$.
\end{theorem}
\begin{proof}\leanok
    \uses{thm:peps_regular_actual_projector,thm:peps_regular_constraint_fixed_spaces,
        thm:peps_regular_flat_union}
    Replace simultaneous vertex and edge invariance in the actual range
    characterization by the corresponding individual fixedness conditions.
    For $R\subseteq S$, every vertex and internal edge of $R$ belongs to $S$,
    and flatness on $S$ implies flatness on $R$.
\end{proof}

\begin{theorem}[Intersection of regular regional ranges]%
    \label{thm:peps_regular_region_intersection}
    \lean{TNLean.PEPS.regularGlobalRegionRange_inf_eq_union}
    \leanok
    \uses{def:peps_ambient_regular_range}
    If $R\cap S=\{b\}$ and no edge joins $R\setminus S$ to $S\setminus R$, then
    \begin{align}
        \widehat{\mathcal S}_R^0\cap\widehat{\mathcal S}_S^0
         & =\widehat{\mathcal S}_{R\cup S}^0.
        \label{eq:peps_regular_region_intersection}
    \end{align}
    The group need not be abelian, and no conditions are imposed on
    complementary half-edge coordinates.
\end{theorem}
\begin{proof}\leanok
    \uses{thm:peps_regular_range_nested,thm:peps_regular_flat_union}
    A vector in both ranges is fixed by all vertex averages on $R\cup S$.
    Every internal edge belongs to one of the two regions, so all required
    edge averages also fix it. A nonzero coefficient must be flat on both
    regions and hence, by Theorem~\ref{thm:peps_regular_flat_union}, on their
    union. The atomic characterization gives membership in the union range.
    The reverse inclusion follows from nesting.
\end{proof}

\begin{definition}[Regional conditions on the core physical space]%
    \label{def:peps_regular_core_slice_range}
    \lean{TNLean.PEPS.regularSubregionSlice,
        TNLean.PEPS.regularRegionRangeWithin,
        TNLean.PEPS.regularRegionConstantExtension}
    \leanok
    \uses{def:peps_ambient_regular_range}
    For $R\subseteq T$, let $\mathcal S_{R;T}^0$ consist of vectors on the
    physical half-edge factors of $T$ whose every slice at fixed coordinates
    on $T\setminus R$ lies in $\mathcal S_R^0$.
    Define $J_Tx$ on the whole graph by
    $(J_Tx)(\alpha)=x(\alpha|_T)$, so it is constant in exterior physical
    coordinates. The space $\mathcal S_{R;T}^0$ itself has physical factors
    only at the vertices of $T$.
\end{definition}

\begin{theorem}[Intersection on only the core physical factors]%
    \label{thm:peps_regular_core_intersection}
    \lean{TNLean.PEPS.regularRegionRangeWithin_eq_comap,
        TNLean.PEPS.regularRegionRangeWithin_self,
        TNLean.PEPS.regularRegionRangeWithin_inf_eq_range}
    \leanok
    \uses{def:peps_regular_core_slice_range}
    One has $\mathcal S_{R;T}^0=J_T^{-1}(\widehat{\mathcal S}_R^0)$ and
    $\mathcal S_{T;T}^0=\mathcal S_T^0=\operatorname{range}\Gamma_T$.
    Suppose $R\cap S=\{b\}$, no edge joins $R\setminus S$ to $S\setminus R$,
    and $T=R\cup S$. Then
    \begin{align}
        \mathcal S_{R;T}^0\cap\mathcal S_{S;T}^0
         & =\operatorname{range}\Gamma_T.
        \label{eq:peps_regular_core_intersection}
    \end{align}
    The vector space has no exterior physical factors, while $\Gamma_T$
    still takes an arbitrary correlated tensor on every boundary virtual leg.
\end{theorem}
\begin{proof}\leanok
    \uses{thm:peps_regular_region_intersection}
    A complementary slice of $J_Tx$ depends only on its labels in
    $T\setminus R$, and every assignment of those labels extends to the
    whole complement by setting exterior labels to the group identity.
    This proves the pullback formula. For $R=T$ each slice is simply $x$.
    Taking the inverse image of (\ref{eq:peps_regular_region_intersection})
    under $J_T$ and using that inverse images preserve intersections gives
    (\ref{eq:peps_regular_core_intersection}).
\end{proof}

\begin{definition}[Three four-legged blocks]%
    \label{def:peps_regular_three_block_geometry}
    \lean{TNLean.PEPS.threeBlockFourLegGraph,
        TNLean.PEPS.threeBlockLeftRegion,TNLean.PEPS.threeBlockRightRegion,
        TNLean.PEPS.threeBlockRegion}
    \leanok
    Take vertices $0,\ldots,10$ and the unordered edges
    \begin{align}
        E & =\{\{0,j\}:j\in\{1,3,4,5\}\}\cup
        \{\{1,j\}:j\in\{2,6,7\}\}\cup
        \{\{2,j\}:j\in\{8,9,10\}\}.
        \label{eq:peps_three_block_edges}
    \end{align}
    Set $R=\{0,1\}$, $S=\{1,2\}$ and $T=\{0,1,2\}$.
    The core is the path $0$--$1$--$2$. Its end vertices each have three
    exterior legs and its middle vertex has two. The eight remaining
    vertices are distinct exterior endpoints, so an arbitrary boundary
    tensor may correlate all eight open legs of $T$.
\end{definition}

The regular open-region map on the three core factors is the following
contraction. Each physical leg carries the four-label coordinate space at
one core vertex; the eight virtual legs carry arbitrary, jointly correlated
boundary data. The two horizontal internal labels are summed.
% Source: SCP10 Theorem 5.4, paper_v3.tex:1373--1403, figs3/grow-intersect.
% This panel is the regular-coordinate specialization, not arbitrary A,B.
% Exact open boundary signature: (virtual,physical)=(8,3).
\begin{tenkzequation}
    \tenkzkernel
    $\Gamma_T^0=$
    % Tenkz lattice dimensions are literal DSL tokens, not printed products.
    \begin{tenkz}[lattice={1x3}, frame=plane, physical=up] % chktex 29
        \tn[at=(1,1), name=intersectionCore0, skin=box,
            ports={180:virtual:$\ell$,0:virtual,90:virtual:$t_0$,270:virtual:$b_0$}, label pos=se]{$\Pi_G$}
        \tn[at=(1,2), name=intersectionCore1, skin=box,
            ports={180:virtual,0:virtual,90:virtual:$t_1$,270:virtual:$b_1$}, label pos=se]{$\Pi_G$}
        \tn[at=(1,3), name=intersectionCore2, skin=box,
            ports={180:virtual,0:virtual:$r$,90:virtual:$t_2$,270:virtual:$b_2$}, label pos=se]{$\Pi_G$}
    \end{tenkz}
\end{tenkzequation}

\begin{theorem}[Four-legged regular intersection]%
    \label{thm:peps_regular_three_block_intersection}
    \lean{TNLean.PEPS.threeBlockFourLegGraph_core_degree,
        TNLean.PEPS.threeBlockFourLegGraph_boundary_card,
        TNLean.PEPS.threeBlock_regularGlobalRegionRange_intersection,
        TNLean.PEPS.threeBlock_regularRegionRange_intersection}
    \leanok
    \uses{def:peps_regular_three_block_geometry,def:peps_regular_core_slice_range}
    Each core vertex has four incident edges, and $T$ has eight distinct
    boundary edges. For every finite group,
    $\widehat{\mathcal S}_R^0\cap\widehat{\mathcal S}_S^0
        =\widehat{\mathcal S}_T^0$ on the full eleven-vertex space, and
    $\mathcal S_{R;T}^0\cap\mathcal S_{S;T}^0
        =\operatorname{range}\Gamma_T$ on exactly the three core physical
    factors. The latter open-region map retains all eight independent
    virtual boundary labels.
\end{theorem}
\begin{proof}\leanok
    \uses{thm:peps_regular_region_intersection,thm:peps_regular_core_intersection}
    Count the edges in (\ref{eq:peps_three_block_edges}). The overlap is
    $\{1\}$, the union is $T$, and the only possible edge between the
    disjoint parts would be $02$, which is absent. Apply
    Theorems~\ref{thm:peps_regular_region_intersection}
    and~\ref{thm:peps_regular_core_intersection}.
\end{proof}
